\begin{center}
  \Large\textbf{KATA PENGANTAR}
\end{center}
\vspace{2ex}

\addcontentsline{toc}{chapter}{KATA PENGANTAR}

% Ubah paragraf-paragraf berikut sesuai dengan yang ingin diisi pada kata pengantar

Puji dan syukur kehadirat Tuhan Yang Maha Esa, yang atas rahmat dan karunia-Nya telah memberikan kelancaran, kekuatan, serta kemudahan, sehingga penulis dapat menyelesaikan magang sekaligus menyusun laporan ini dengan baik dan tepat waktu.

Penyusunan laporan ini tentu tidak akan terwujud tanpa adanya bimbingan, dukungan, dan semangat dari berbagai pihak. Oleh karena itu, pada kesempatan ini penulis ingin menyampaikan ucapan terima kasih yang tulus kepada:


\begin{enumerate}[nolistsep]

  \item Bapak Dr. Arief Kurniawan, S.T., M.T. selaku Kepala Departemen Teknik Komputer FTEIC-ITS sekaligus Dosen Pembimbing Magang penulis.

  \item Ibu Dr. Diah Puspito Wulandari, S.T., M.Sc. selaku Koordinator Magang Departemen Teknik Komputer FTEIC-ITS.

  \item Bapak Edi Gunawan, S.Kom. selaku Pembimbing Lapangan di Badan Pendapatan Daerah Kota Surabaya selama pelaksanaan magang.
  
  \item Rekan-rekan di Badan Pendapatan Daerah Kota Surabaya yang telah memberikan bantuan, dukungan, serta ilmu pengetahuan selama pelaksanaan magang.
  
  \item Semua pihak yang tidak dapat penulis sebutkan satu per satu, yang telah membantu dan mendukung penulis hingga penyusunan laporan ini.

\end{enumerate}

Akhir kata, semoga laporan ini dapat memberikan manfaat dan wawasan bagi semua pihak yang membacanya. Penulis sadar bahwa laporan ini masih jauh dari kata sempurna sehingga penulis menerima segala saran dan kritik yang membangun. Penulis berharap laporan ini dapat memberikan kontribusi dalam perkembangan ilmu pengetahuan di bidang teknologi informasi.

\begin{flushright}
  \begin{tabular}[b]{c}
    % Ubah kalimat berikut sesuai dengan tempat, bulan, dan tahun penulisan
    Surabaya, Desember 2025
    \\
    \\
    \\
    \\
    Penulis
  \end{tabular}
\end{flushright}
